\section{Conclusion}\label{sec:conclusion}

We asked whether one spatio-temporal model can jointly forecast next-day regional evening-peak demand and a graph-level Operational Stress Index from Bangladesh national grid records under chronological evaluation. The answer, on this dataset and protocol, is yes. Correlation-graph PF-STGT with parallel spatial and temporal encoding and repaired multi-task loss gives the best joint results in our benchmark suite: macro demand MAE 88.65~MW and OSI $R^2$ 0.745, with statistically significant demand improvement over the historical hybrid reference.

Three findings stand out. Correlation adjacency works better than geographical priors here. Joint training trades a small amount of demand accuracy for substantially more useful stress output compared with demand-only training. Attribution analysis points to limitation and calendar coalitions for stress and to calendar, lag, and generation channels for Dhaka demand, with only partial agreement across methods.

The study adds a reproducible Bangladesh case for coupled load-and-stress forecasting and shows that graph design and loss weighting need empirical tuning rather than assumptions alone. A single model can support paired day-ahead forecasts when graph structure, multi-task calibration, and division-level monitoring are designed together---with the usual caveat that our evidence is from offline evaluation, not from demonstrated field deployment.
